\documentclass[a4paper,11pt]{article}
        \usepackage[top=15mm,bottom=18mm,left=16mm,right=16mm]{geometry}
        \usepackage{fontspec}
        \usepackage{xeCJK}
        \setmainfont{Times New Roman}
        \setCJKmainfont{Yu Gothic}
        \setmonofont{Consolas}
        \setCJKmonofont{Yu Gothic}
        \usepackage{booktabs}
        \usepackage{longtable}
        \usepackage{tabularx}
        \usepackage{array}
        \usepackage{enumitem}
        \usepackage{hyperref}
        \usepackage{listings}
        \usepackage{xcolor}
        \hypersetup{
          colorlinks=true,
          linkcolor=blue,
          urlcolor=blue,
          pdftitle={dashboard + HTTP API},
          pdfauthor={Codex}
        }
        \lstset{
          basicstyle=\ttfamily\small,
          breaklines=true,
          columns=fullflexible,
          keepspaces=true,
          frame=single,
          framerule=0.2pt,
          rulecolor=\color{black},
          showstringspaces=false
        }
        \setlength{\parskip}{0.42em}
        \setlength{\parindent}{1em}
        \renewcommand{\arraystretch}{1.15}
        \title{30. dashboard + HTTP API}
        \author{solar\_ws0129-main}
        \date{2026-07-29}
        \begin{document}
        \maketitle

        \section*{対象}
        \begin{tabularx}{\linewidth}{>{\raggedright\arraybackslash}p{0.18\linewidth} X}
        \toprule
        項目 & 内容 \\
        \midrule
        ファイル & \path|mpc_solarcar/dashboard_node.py| \\
        種別 & Python \\
        区分 & runtime node \\
        役割 & planner と vehicle の現在値を ROS から受け、HTTP サーバと dashboard frontend へ渡す。 \\
        起動文脈 & 可視化の中心ノード。 \\
        \bottomrule
        \end{tabularx}

        \section*{このファイルを読む理由}
        planner と vehicle の現在値を ROS から受け、HTTP サーバと dashboard frontend へ渡す。

        \begin{itemize}[leftmargin=1.6em]
        \item /api/state と /metrics を持つ。
\item ROS topic を直接 browser へ出すのではなく、Node 内 state に集約する。
        \end{itemize}

        \section*{呼び出し関係}
        \subsection*{どこから呼ばれるか}
        \begin{itemize}[leftmargin=1.6em]
        \item \path|launch/solarcar_sim.launch.py|
\item \path|mpc_solarcar/live_launch.py|
\item \path|launch/solar_measurement.launch.py|
        \end{itemize}

        \subsection*{次に読むと理解が進むファイル}
        \begin{itemize}[leftmargin=1.6em]
        \item 特記事項なし。
        \end{itemize}

        \subsection*{同一リポジトリ内の主要依存}
        \begin{itemize}[leftmargin=1.6em]
        \item 同一リポジトリ内の明示 import は少ない。
        \end{itemize}

        \section*{ソース冒頭の把握}
        モジュール docstring または先頭意図の要約:

        モジュール docstring は明示されていない。

        \subsection*{代表コード断片}
        \begin{lstlisting}[language=Python]
from ament_index_python.packages import get_package_share_directory


class DashboardHandler(SimpleHTTPRequestHandler):
    def __init__(self, *args, node=None, directory=None, **kwargs):
        self._node = node
        super().__init__(*args, directory=directory, **kwargs)

    def do_GET(self):
        if self.path.startswith('/api/state'):
            self._send_json(self._node.get_state())
            return
        if self.path.startswith('/metrics'):
            self._send_metrics(self._node.get_prometheus_metrics())
            return
        super().do_GET()

    def log_message(self, format, *args):
        # quiet
        return

    def _send_json(self, data):
\end{lstlisting}


        \section*{主要構造の読み取り}
        主要クラスは DashboardHandler, DashboardNode。 主要関数は do\_GET, log\_message, destroy\_node, get\_state, get\_prometheus\_metrics, main。 ROS パラメータ宣言は 5 件。 ROS I/O は publisher 0 件、subscription 21 件。

        \subsection*{主要クラス・関数}
            \begin{longtable}{p{0.07\linewidth} p{0.16\linewidth} p{0.25\linewidth} p{0.40\linewidth}}
            \toprule
            行 & 種別 & 名前 & 補足 \\
            \midrule
            19 & class & \path|DashboardHandler| & bases: SimpleHTTPRequestHandler \\
55 & class & \path|DashboardNode| & bases: Node \\
20 & function & \path|__init__| & args: self \\
24 & function & \path|do_GET| & args: self \\
33 & function & \path|log_message| & args: self, format \\
37 & function & \path|_send_json| & args: self, data \\
45 & function & \path|_send_metrics| & args: self, payload \\
56 & function & \path|__init__| & args: self \\
156 & function & \path|destroy_node| & args: self \\
165 & function & \path|get_state| & args: self \\
169 & function & \path|get_prometheus_metrics| & args: self \\
236 & function & \path|_update| & args: self, key, value \\
241 & function & \path|_on_speed_cmd| & args: self, msg \\
244 & function & \path|_on_upper_speed_cmd| & args: self, msg \\
247 & function & \path|_on_speed_meas| & args: self, msg \\
250 & function & \path|_on_throttle_cmd| & args: self, msg \\
253 & function & \path|_on_drive_mode| & args: self, msg \\
256 & function & \path|_on_soc| & args: self, msg \\
264 & function & \path|_set_estimated_soc| & args: self, value \\
273 & function & \path|_on_tb| & args: self, msg \\
276 & function & \path|_on_ibatt| & args: self, msg \\
279 & function & \path|_on_vbatt| & args: self, msg \\
282 & function & \path|_on_s_km| & args: self, msg \\
285 & function & \path|_on_status| & args: self, msg \\
301 & function & \path|_on_env| & args: self, msg \\
313 & function & \path|_on_metrics| & args: self, msg \\
333 & function & \path|_on_upper_plan| & args: self, msg \\
344 & function & \path|_on_lower_plan| & args: self, msg \\
355 & function & \path|_on_sys_state| & args: self, msg \\
358 & function & \path|_on_sys_diag| & args: self, msg \\
361 & function & \path|_on_mpc_state| & args: self, msg \\
364 & function & \path|_on_health| & args: self, msg \\
367 & function & \path|_on_gps| & args: self, msg \\
373 & function & \path|_on_path| & args: self, msg \\
379 & function & \path|_init_dummy| & args: self, path, rate\_hz \\
395 & function & \path|_tick_dummy| & args: self \\
409 & function & \path|main| &  \\
            \bottomrule
            \end{longtable}

        \subsection*{ファイルを上から読んだときの定義順}
            \begin{longtable}{p{0.07\linewidth} p{0.84\linewidth}}
            \toprule
            行 & 内容 \\
            \midrule
            19 & クラス DashboardHandler を定義する。 \\
55 & クラス DashboardNode を定義する。 \\
409 & 関数 main を定義する。 \\
417 & 条件 \_\_name\_\_ == '\_\_main\_\_' を判定し、真なら内部処理を行う。 \\
            \bottomrule
            \end{longtable}

        \subsection*{import 群の詳細}
            \begin{longtable}{p{0.07\linewidth} p{0.33\linewidth} p{0.50\linewidth}}
            \toprule
            行 & import 文 & このファイルで読む理由 \\
            \midrule
            1 & \path|import json| & manifest、checkpoint、UDP payload をやり取りするため。 このファイル内での主な使用位置は L38。 \\
2 & \path|import math| & clip、sqrt、sin/cos、有限判定など数値ロジックに使うため。 このファイル内での主な使用位置は L212。 \\
3 & \path|import os| & パス、環境変数、プロセス外部状態を扱うため。 このファイル内での主な使用位置は L68, L70, L71。 \\
4 & \path|import threading| & threading モジュールを利用するため。 このファイル内での主な使用位置は L73, L150。 \\
5 & \path|import time| & wall/monotonic time に基づく周期制御や freshness 判定を行うため。 このファイル内での主な使用位置は L75, L215, L239, L261, L262, L268, L299, L311, ...。 \\
6 & \path|from functools import partial| & functools から partial を読み込み、このファイルの処理を組み立てるため。 このファイル内での主な使用位置は L147。 \\
7 & \path|from http.server import SimpleHTTPRequestHandler, ThreadingHTTPServer| & http.server から SimpleHTTPRequestHandler, ThreadingHTTPServer を読み込み、このファイルの処理を組み立てるため。 このファイル内での主な使用位置は L19, L149。 \\
9 & \path|import rclpy| & ROS 2 Python ノードとして起動・spin するため。 このファイル内での主な使用位置は L410, L412, L414。 \\
10 & \path|from rclpy.node import Node| & ROS 2 ノード本体の基底クラスとして使うため。 このファイル内での主な使用位置は L55。 \\
12 & \path|from std_msgs.msg import Float32, Float32MultiArray, String| & Float32 や String など軽量メッセージで planner / vehicle 値を publish するため。 このファイル内での主な使用位置は L117, L118, L119, L120, L121, L122, L123, L124, ...。 \\
13 & \path|from sensor_msgs.msg import NavSatFix| & GPS 緯度経度高度を ROS topic でやり取りするため。 このファイル内での主な使用位置は L136, L367。 \\
14 & \path|from nav_msgs.msg import Path| & 将来軌跡を dashboard や logger へ出すため。 このファイル内での主な使用位置は L137, L373。 \\
16 & \path|from ament_index_python.packages import get_package_share_directory| & ament\_index\_python.packages から get\_package\_share\_directory を読み込み、このファイルの処理を組み立てるため。 このファイル内での主な使用位置は L67。 \\
381 & \path|import pandas as pd| & CSV や時刻列の読込・整形・再サンプリングに使うため。 このファイル内での主な使用位置は L382。 \\
            \bottomrule
            \end{longtable}

        \section*{関数・クラスを上から順に解説}
        \subsection*{L19 クラス \path|DashboardHandler|}
            \begin{itemize}[leftmargin=1.6em]
              \item 定義: \path|DashboardHandler(bases=SimpleHTTPRequestHandler)|
              \item docstring: 明示されていない。
              \item このブロックが直接呼ぶ主な関数・メソッド: \_\_init\_\_, \_send\_json, \_send\_metrics, do\_GET, dumps, encode, end\_headers, get\_prometheus\_metrics, get\_state, len, send\_header, send\_response
              \item 戻り値の要点: 戻り値が無いか、return 文が明示されていない。
            \end{itemize}
            \begin{enumerate}[leftmargin=1.8em]
            \item 関数 \_\_init\_\_ を定義する。
\item 関数 do\_GET を定義する。
\item 関数 log\_message を定義する。
\item 関数 \_send\_json を定義する。
\item 関数 \_send\_metrics を定義する。
            \end{enumerate}

\subsection*{L55 クラス \path|DashboardNode|}
            \begin{itemize}[leftmargin=1.6em]
              \item 定義: \path|DashboardNode(bases=Node)|
              \item docstring: 明示されていない。
              \item このブロックが直接呼ぶ主な関数・メソッド: Lock, Thread, ThreadingHTTPServer, \_\_init\_\_, \_init\_dummy, \_set\_estimated\_soc, \_update, abspath, create\_subscription, create\_timer, declare\_parameter, destroy\_node
              \item 戻り値の要点: 戻り値が無いか、return 文が明示されていない。
            \end{itemize}
            \begin{enumerate}[leftmargin=1.8em]
            \item 関数 \_\_init\_\_ を定義する。
\item 関数 destroy\_node を定義する。
\item 関数 get\_state を定義する。
\item 関数 get\_prometheus\_metrics を定義する。
\item 関数 \_update を定義する。
\item 関数 \_on\_speed\_cmd を定義する。
\item 関数 \_on\_upper\_speed\_cmd を定義する。
\item 関数 \_on\_speed\_meas を定義する。
\item 関数 \_on\_throttle\_cmd を定義する。
\item 関数 \_on\_drive\_mode を定義する。
\item 関数 \_on\_soc を定義する。
\item 関数 \_set\_estimated\_soc を定義する。
            \end{enumerate}

\subsection*{L409 関数 \path|main|}
            \begin{itemize}[leftmargin=1.6em]
              \item 定義: \path|main()|
              \item docstring: 明示されていない。
              \item このブロックが直接呼ぶ主な関数・メソッド: DashboardNode, destroy\_node, init, shutdown, spin
              \item 戻り値の要点: 戻り値が無いか、return 文が明示されていない。
            \end{itemize}
            \begin{enumerate}[leftmargin=1.8em]
            \item rclpy.init(...) を実行する。
\item node に DashboardNode() の結果を代入する。
\item rclpy.spin(...) を実行する。
\item node.destroy\_node(...) を実行する。
\item rclpy.shutdown(...) を実行する。
            \end{enumerate}

        \subsection*{設定値と入力パラメータ}
            \begin{longtable}{p{0.07\linewidth} p{0.35\linewidth} p{0.45\linewidth}}
            \toprule
            行 & 名前 & 既定値・補足 \\
            \midrule
            58 & \path|host| & 0.0.0.0 \\
59 & \path|port| & 8080 \\
60 & \path|static_dir| &  \\
61 & \path|dummy_csv| &  \\
62 & \path|dummy_rate_hz| & 5.0 \\
            \bottomrule
            \end{longtable}

        \subsection*{ROS topic I/O}
            \begin{longtable}{p{0.07\linewidth} p{0.18\linewidth} p{0.34\linewidth} p{0.28\linewidth}}
            \toprule
            行 & 種別 & topic & callback / 補足 \\
            \midrule
            117 & subscription & \path|/planner/speed_cmd| & self.\_on\_speed\_cmd \\
118 & subscription & \path|/planner/upper_speed_cmd| & self.\_on\_upper\_speed\_cmd \\
119 & subscription & \path|/vehicle/speed_kmh| & self.\_on\_speed\_meas \\
120 & subscription & \path|/planner/throttle_cmd_pct| & self.\_on\_throttle\_cmd \\
121 & subscription & \path|/planner/drive_mode| & self.\_on\_drive\_mode \\
122 & subscription & \path|/vehicle/batt_soc| & self.\_on\_soc \\
123 & subscription & \path|/vehicle/batt_temp_c| & self.\_on\_tb \\
124 & subscription & \path|/vehicle/batt_current_a| & self.\_on\_ibatt \\
125 & subscription & \path|/vehicle/batt_voltage_v| & self.\_on\_vbatt \\
126 & subscription & \path|/vehicle/s_km| & self.\_on\_s\_km \\
127 & subscription & \path|/planner/status| & self.\_on\_status \\
128 & subscription & \path|/planner/env| & self.\_on\_env \\
129 & subscription & \path|/planner/metrics| & self.\_on\_metrics \\
130 & subscription & \path|/planner/upper_plan| & self.\_on\_upper\_plan \\
131 & subscription & \path|/planner/lower_plan| & self.\_on\_lower\_plan \\
132 & subscription & \path|/system/state| & self.\_on\_sys\_state \\
133 & subscription & \path|/system/diag| & self.\_on\_sys\_diag \\
134 & subscription & \path|/system/mpc_state| & self.\_on\_mpc\_state \\
135 & subscription & \path|/system/health| & self.\_on\_health \\
136 & subscription & \path|/sim/gps| & self.\_on\_gps \\
137 & subscription & \path|/planner/trajectory| & self.\_on\_path \\
            \bottomrule
            \end{longtable}







        \section*{処理の流れ}
        \begin{enumerate}[leftmargin=1.7em]
        \item 初期化時に設定値や入力パスを読み込む。
\item publisher / subscription / timer を準備する。
\item timer callback 周期で主処理を進める。
        \end{enumerate}

        \section*{このファイルの位置づけ}
        この資料群では、\path|mpc_solarcar/dashboard_node.py| を
        \textbf{runtime node}の中核として扱う。
        したがって、ここで扱う処理は単独で完結せず、
        前段の profile / launch / telemetry と後段の planner / logger / report へ繋がる。

        \end{document}
